\setcounter{ex}{0}
\setcounter{paragraph}{0}
\PhanI
\begin{ex}
Chuyển động nào sau đây \textbf{không} phải là dao động cơ học?
\choice
{Chuyển động đung đưa của con lắc của đồng hồ}
{Chuyển động đung đưa của lá cây}
{Chuyển động nhấp nhô của phao trên mặt nước}
{\True Chuyển động thẳng đều của ôtô trên đường}
\loigiai{}
\end{ex}
\begin{ex} %[3P1Y1-1][Loai1] 
(QG - 2018) Một vật dao động điều hòa theo phương trình $x = A\cos\left(\omega t + \varphi \right)$ ($\omega  > 0$). Tần số góc của dao động là
\choice
{A}
{\True $\omega$}
{$\varphi$}
{x}
\loigiai{
Tần số góc của dao động là $\omega$.
}
\end{ex} \haidong{12}
\begin{ex} %[3P1Y1-1][Loai1] 
(QG - 2018) Một vật dao động điều hòa theo phương trình $x = A\cos\left(\omega t + \varphi \right)$ ($A > 0$). Biên độ dao động của vật là
\choice
{\True A}
{$\varphi$}
{$\omega$}
{x}
\loigiai{
Biên độ dao động của vật là A.
}
\end{ex} \haidong{12}
\begin{ex} %[3P1Y1-1][Loai1] 
(MH - 2017) Một chất điểm dao động điều hòa với phương trình $x=A\cos\left(\omega t+\varphi\right)$, trong đó $\omega $ có giá trị dương. Đại lượng $\omega$ gọi là
\choice
{biên độ dao động}
{chu kì của dao động}
{\True tần số góc của dao động}
{pha ban đầu của dao động}
\loigiai{
$\omega$ được gọi là tần số góc của dao động.
}
\end{ex} \haidong{12}
\begin{ex} %[3P1Y1-1][Loai1] 
(MH - 2018) Một vật dao động điều hòa trên trục Ox quanh vị trí cân bằng O. Gọi $A,\ \omega$ và $\varphi$ lần lượt là biên độ, tần số góc và pha ban đầu của dao động. Biểu thức li độ của vật theo thời gian t là
\choice
{\True $x = A\cos\left(\omega t + \varphi\right)$}
{$x = \omega \cos\left(t\varphi +A\right) $}
{$x = t\cos\left(\varphi A + \omega\right)$}
{$x = \varphi \cos\left(A\omega +t\right)$}
\loigiai{
Biểu thức li độ của vật theo thời gian là $x = A\cos\left(\omega t + \varphi\right)$.
}
\end{ex} \haidong{12}
\begin{ex} %[3P1Y1-1][Loai1] 
(MH - 2019) Một vật dao động điều hoà theo phương trình $x = A\cos\left(\omega t + \varphi\right)$ ($A >0$, $\omega > 0$). Pha của dao động ở thời điểm t là
\choice
{$\omega$}
{$\cos\left(\omega t + \varphi\right)$}
{\True $\omega t + \varphi$}
{$\varphi$}
\loigiai{
Pha của dao động $\alpha =\omega t+\varphi $.
}
\end{ex} \haidong{12}
\begin{ex}%[3P1Y1-1][Loai1]
Một vật dao động điều hòa với biên độ A. Giá trị cực tiểu của li độ trong quá trình vật dao động là
\choice
{A}
{0}
{\True $-A$}
{$-2A$}
\loigiai{
Giá trị cực tiểu của li độ trong quá trình vật dao động là $-A$.
}
\end{ex} \haidong{12}
% 
\begin{ex}
Chu kì dao động điều hòa là
\choice
{số dao động toàn phần vật thực hiện được trong $1 \mathrm{~s}$}
{khoảng thời gian để vật đi từ bên này sang bên kia của quỹ đạo chuyển động}
{khoảng thời gian ngắn nhất để vật trở lại vị trí ban đầu}
{\True khoảng thời gian ngắn nhất để vật lặp lại trạng thái dao động}
\end{ex} \haidong{12}

\begin{ex} %[3P1B1-1][Loai1]  
Phát biểu nào sau đây là \textbf{sai}? Dao động cơ điều hòa là một
\choice
{chuyển động có giới hạn trong không gian}
{dao động cơ học}
{dao động tuần hoàn}
{\True chuyển động có quỹ đạo hình sin}
\loigiai{
\begin{itemize}
\item Dao động điều hòa là một chuyển động lặp đi lặp lại quanh một vị trí cân bằng, mà li độ dao động là một hàm sin của thời gian nên D sai.
\end{itemize}}
\end{ex} \haidong{12}

\begin{ex}
Tần số dao động điều hòa là
\choice
{\True số dao động toàn phần vật thực hiện được trong $1 \mathrm{~s}$}
{số dao động toàn phần vật thực hiện được trong một chu kì}
{khoảng thời gian ngắn nhất để vật trở lại vị trí ban đầu}
{khoảng thời gian vật thực hiện hết một dao động toàn phần}
\end{ex} \haidong{12}

\begin{ex}
Phát biểu nào sau đây đúng về dao động điều hoà?
\choice
{Quỹ đạo là đường hình sin}
{\True Quỹ đạo là một đoạn thẳng}
{Vận tốc tỉ lệ thuận với thời gian}
{Gia tốc tỉ lệ thuận với thời gian}
\end{ex} \haidong{12}
\begin{ex}
Pha dao động của một chất điểm dao động điều hòa dùng để xác định
\choice
{chu kì dao động}
{tần số dao động}
{biên độ dao động}
{\True trạng thái dao động}
\end{ex} \haidong{12}
\begin{ex} %[3P1Y1-1][Loai1]  
Khi một chất điểm dao động điều hòa thì li độ của chất điểm là
\choice
{\True một hàm sin của thời gian}
{là một hàm tan của thời gian}
{là một hàm bậc nhất của thời gian}
{là một hàm bậc hai của thời gian}
\loigiai{
\begin{itemize}
\item Dao động điều hòa là một dao động trong đó li độ của vật là một hàm côsin (hay sin) của thời gian.
\end{itemize}}
\end{ex} \haidong{12}
\begin{ex}%[3P1B1-1]
Một vật dao động điều hòa thì pha của dao động
\choice
{biến thiên điều hòa theo thời gian}
{\True là hàm bậc nhất của thời gian}
{là hàm bậc hai của thời gian}
{không đổi theo thời gian}
\loigiai{
\begin{itemize}
\item Pha dao động tại thời điểm t: $\phi_t=\omega t+\varphi$ là hàm bậc nhất của thời gian t.
\end{itemize}
}
\end{ex} \haidong{12}
\begin{ex}
Một vật dao động điều hòa với phương trình li độ $x=A \cos (\omega t+\varphi)$, đại lượng luôn tăng đều theo thời gian là
\choice
{\True pha dao động $\omega t+\varphi$}
{li độ $x$}
{biên độ $A$}
{pha ban đầu $\varphi$}
\loigiai{
Trong phương trình dao động điều hòa $x=A \cos (\omega t+\varphi)$, pha dao động $\omega t+\varphi$ là một hàm bậc nhất của thời gian $t$ với hệ số góc $\omega > 0$, do đó nó tăng đều theo thời gian.
}
\end{ex}
\begin{ex}
Một vật dao động điều hòa theo phương trình $x=4 \cos \left(4 \pi t+\dfrac{\pi}{3}\right)\ (cm)$ (t tính bằng giây). Pha ban đầu của dao động là
\choice
{$4\pi\ rad$}
{\True $\dfrac{\pi}{3}\ rad$}
{$4\ rad$}
{$\left(4\pi t+\dfrac{\pi}{3}\right)\ rad$}
\loigiai{
Phương trình dao động điều hòa có dạng $x = A \cos(\omega t + \varphi)$, trong đó $\varphi$ là pha ban đầu. So sánh với phương trình đã cho, ta thấy pha ban đầu là $\dfrac{\pi}{3} rad$.
}
\end{ex}
\begin{ex}
\immini[thm]{Một chất điểm dao động điều hòa có li độ phụ thuộc thời gian theo hàm cosin như hình vẽ. Chất điểm có biên độ là 
\choice
{\True $4\ cm$}
{$2\ cm$}
{$-4\ cm$}
{$-2\ cm$}
}{
\begin{hinhanh}{7}
\begin{tikzpicture}[very thick,>=stealth']
\begin{axis}[
width=7cm, height=4.5cm,
axis lines=middle,
axis line style={very thick,Mapcolor}, % <--- dòng thêm vào
xmin=0, xmax=1.8*pi,
ymin=-4.5, ymax=5,
xlabel style={above},
ylabel style={above},
grid=major,
x grid style={dashed, Mapcolor},
y grid style={dashed, Mapcolor},
clip=true,
enlargelimits=false,
extra y ticks={0},
extra y tick label={0},
xlabel={$t~(\mathrm{ms})$},
ylabel={$x~(\mathrm{cm})$},
xtick={0, 1/2, 1, 3/2, 2,2.5,3,3.5,4,4.5,5},
xticklabels={,,,,20,},
ytick={-4, -2, 0, 2, 4},
yticklabels={-4, -2, 0, 2, 4},
]
\addplot[line width=1.2pt, domain=0:1.5*pi, samples=400, Mapcolor]{4*cos(deg(0.5*pi*x))};
\addplot[Mapcolor,mark=*,mark options={fill=white,mark size=1.5pt}, nodes near coords={}]
coordinates {(2, 0)};
\end{axis}
\end{tikzpicture}
\end{hinhanh}}
\loigiai{}
\end{ex}
\begin{ex}
Một chất điểm dao động điều hòa có đồ thị li độ theo thời gian như hình vẽ. Chu kì dao động là\\
\centerline{\begin{tikzpicture}[very thick,>=stealth']
\begin{axis}[
width=9cm, height=4.5cm,
axis lines=middle,
axis line style={very thick,Mapcolor}, % <--- dòng thêm vào
xmin=0, xmax=3.8*pi,
ymin=-4.5, ymax=5,
xlabel style={above},
ylabel style={above},
grid=major,
x grid style={dashed, Mapcolor},
y grid style={dashed, Mapcolor},
clip=true,
enlargelimits=false,
extra y ticks={0},
extra y tick label={0},
xlabel={$t~(\mathrm{s})$},
ylabel={$x~(\mathrm{cm})$},
xtick={1,2,3,4,5,6,7,8,9,10,11,12},
xticklabels={,,,,,,,0.8},
ytick={-4, -2, 0, 2, 4},
yticklabels={},
]
\addplot[line width=1.2pt, domain=0:3.5*pi, samples=400, Mapcolor]{4*cos(deg(0.5*pi*x))};
\addplot[Mapcolor,mark=*,mark options={fill=white,mark size=1.5pt}, nodes near coords={}]
coordinates {(8, 0)};
\end{axis}
\end{tikzpicture}}
\choice
{$0,8\ s$}
{$0,1\ s$}
{$0,2\ s$}
{\True $0,4\ s$}
\loigiai{}
\end{ex}
\setcounter{ex}{0}
\PhanII
\begin{ex}
Dao động điều hòa là dạng chuyển động cơ học cơ bản, trong đó vị trí của vật được mô tả bởi các hàm $\sin$ hoặc $\cos$ theo thời gian.
\choiceTF[t]
{Các đại lượng $A,\ \omega,\ \varphi$ dao động điều hòa cùng tần số}
{\True Pha dao động là hàm bậc nhất theo thời gian}
{\True Biên độ của vật luôn dương}
{Quỹ đạo dao động của vật là một đường thẳng}
\loigiai{
\begin{itemchoice}
\itemch Sai: $A,\ \omega,\ \varphi$ là các hằng số, không dao động.
\itemch Đúng: Pha là $\omega t + \varphi$.
\itemch Đúng: Biên độ là độ dịch chuyển cực đại, luôn dương.
\itemch Dao động điều hòa xảy ra trên một đoạn thẳng.
\end{itemchoice}
}
\end{ex}
\begin{ex}
Dao động điều hòa thường xuất hiện ở các hệ cơ học đơn giản như con lắc lò xo hay con lắc đơn.
\choiceTF[t]
{\True Dao động tuần hoàn đơn giản nhất là dao động điều hòa}
{\True Pha dao động giúp ta xác định trạng thái dao động của vật}
{Biên độ là độ dịch chuyển của vật tính từ vị trí cân bằng}
{\True Chu kì là khoảng thời gian ngắn nhất để vật lặp lại một trạng thái dao động}
\loigiai{
\begin{itemchoice}
\itemch Đúng: Đây là định nghĩa cơ bản.
\itemch Đúng: Pha xác định vị trí và chiều chuyển động của vật tại thời điểm bất kỳ.
\itemch Sai: Biên độ là độ dịch chuyển cực đại, không phải bất kỳ.
\itemch Đúng: Chu kì là thời gian nhỏ nhất để trạng thái dao động lặp lại.
\end{itemchoice}
}
\end{ex}
\begin{ex}
Dao động điều hòa là chuyển động nền tảng của nhiều hệ cơ học (con lắc lò xo, con lắc đơn …).
\choiceTF[t]
{Khoảng thời gian ngắn nhất sau đó trạng thái dao động lặp lại như cũ gọi là tần số dao động}
{\True Đại lượng không đổi theo thời gian là biên độ dao động}
{Quỹ đạo dao động của vật biểu diễn dưới dạng hình sin}
{Pha ban đầu của dao động luôn dương}
\loigiai{
\begin{itemchoice}
\itemch
\itemch
\itemch
\itemch
\end{itemchoice}
}
\end{ex}
\haidong{12}




\begin{ex}
Ở những hệ quen thuộc như con lắc đồng hồ hay lò xo của hệ thống giảm xóc của xe máy, vật dao động qua lại quanh vị trí cân bằng.
\choiceTF[t]
{\True Dao động cơ học là chuyển động qua lại của một vật quanh vị trí cân bằng}
{Dao động cơ của một vật là dao động tuần hoàn}
{\True Dao động điều hòa là dao động trong đó li độ được mô tả bằng một định luật dạng cosin (hay sin) theo thời gian}
{\True Dao động điều hòa là trường hợp đặc biệt của dao động tuần hoàn}
\loigiai{
\begin{itemchoice}
\itemch
\itemch
\itemch
\itemch
\end{itemchoice}
}
\end{ex} \haidong{12}
